\medskip


\begin{center}A(0~;~4~;~16), \quad B$(0~;~4~;~- 10)$, \quad C$(4~;~-8~;~0)$\quad et\quad  K(0~;~4~;~3).\end{center}




\medskip

\begin{enumerate}
\item 
	\begin{enumerate}
		\item 
		On a CM$^2 = (4 - 0)^2 + (- 8 - 4)^2 + (0 - 3)^2 = 16 + 144 + 9 = 169 = 13^2$, donc 
		
CM $ = 13 \iff \text{C} \in S$.
		\item 
Calculons les coordonnées de $\vect{\text{AC}}\begin{pmatrix}4\\-12\\-16\end{pmatrix}, \: \vect{\text{BC}}\begin{pmatrix}4\\-12\\10\end{pmatrix}$

D'où le produit scalaire : $\vect{\text{AC}} \cdot \vect{\text{BC}} = 16 + 144 - 160 = 160 - 160 = 0$ : les vecteurs sont orthogonaux, les droites (AC) et (BC) sont perpendiculaires en C, donc le triangle ABC est rectangle en C.
	\end{enumerate}
\item 
	\begin{enumerate}
		\item ~
$\bullet~~$$\vect{n} \cdot \vect{\text{AC}} = 12 - 12 + 0 = 0$ : $\vect{n}$ est orthogonal à $\vect{\text{AC}}$ ;

$\bullet~~$$\vect{n} \cdot \vect{\text{BC}} = 12 - 12 + 0 = 0$ : $\vect{n}$ est orthogonal à $\vect{\text{BC}}$.

Les vecteurs $\vect{\text{AC}}$ et $\vect{\text{BC}}$ ne sont pas colinéaires (leurs coordonnées sont égales mais pas les dernières), donc $\vect{n}$ normal deux vecteurs non colinéaires du plan (ABC) est normal à ce plan.
		\item 
		On sait qu'alors les coordonnées de $\vect{n}$ sont les coefficients de $x$, \: $y$ et $z$ dans l'équation du plan (ABC).
		
On a donc $M(x~;~y~;~z) \in (\text{ABC}) \iff 3x + y + d = 0$, avec $d \in \R$.
		
Or par exemple :
		
C$(4~;~-8~;~0) \in (\text{ABC}) \iff 3\times 4  + (- 8) + d = 0 \iff 12 - 8 + d = 0 \iff d = - 4$.
		
Finalement $M(x~;~y~;~z) \in (\text{ABC}) \iff 3x + y - 4 = 0$.
	\end{enumerate}
\item 


	\begin{enumerate}
		\item 
Si D appartient à l'axe des abscisses son ordonnée et sa cote sont nuls, donc D a pour coordonnées $(x~;~0~;~;0)$.

De plus KD $ = 13 \Rightarrow \text{KD}^2 = 13^2 = 169 = (x - 0)^2 + (0 - 4)^2 + (0 - 3)^2\iff (x - 0)^2 + 16 + 9 = 169 \iff \iff x^2 = 144 \iff x^2 = 12^2 \iff x ^2 - 12^2 = 0 \iff (x + 12)(x - 12) = 0\left\{\begin{array}{l c l}
x - 4 + 12&=&0\\
x - -&=&0
\end{array}\right. \iff \left\{\begin{array}{l c l}
x + 12&=&0\\
x - 12&=&0
\end{array}\right. \iff \left\{\begin{array}{l c l}
x &=&- 12\\
x &=& 12
\end{array}\right.$

Conclusion D(12~;~0~;~0).
		\item 
Si la droite $\Delta$ est perpendiculaire au plan (ABC) un de ses vecteurs directeurs est le vecteur $\vect{n}$ normal à ce plan. On a donc :

$M(x~;~y~;~z) \in \Delta \iff \vect{\text{D}M} = t\vect{n}$, avec $t \in \R$. Cette égalité se traduit par le système :

$\left\{\begin{array}{l c l}
x - 12&=&3t\\
y - 0&=&1t\\
z - 0&=&0t
\end{array}\right. \iff \left\{\begin{array}{l c r}
x &=&12 + 3t\\
y &=&t\\
z &=&0t
\end{array}\right.$.
		\item  
		D appartient à $\Delta$ perpendiculaire au plan (ABC) ; cherchons les coordonnées du projeté orthogonal de D sur le plan (ABC) : I
		
		I appartient au plan (ABC) et appartient à $\Delta$, ses coordonnées $(x~;~y~;~z)$ vérifient donc l'équation du plan (ABC) et les équations paramétriques de $\Delta$ soit le système :
		
$\left\{\begin{array}{l c l}
x &=&12 + 3t\\
y &=&t\\
z &=&0t\\
3x + y - 4&=&0
\end{array}\right.$

En remplaçant dans la dernière équation $x$ et $y$ par leurs valeurs en fonction de $t$ on obtient $3(12 + 3t) + t - 4 = 0 \iff 36 + 9t + t - 4 = 0 \iff 10t + 32 = 0$

$\iff 10t = - 32 \iff t = - \dfrac{32}{10} = - \dfrac{16}{5}$.

En reportant cette valeur dans les deux premières équations de $\Delta$, on obtient :

$x = 12 -3 \times \dfrac{16}{5} = \dfrac{60}{5} - \dfrac{48}{5} = \dfrac{12}{5}$ ;
$y = - \dfrac{16}{5}$ et $z = 0$.

I$\left(\dfrac{12}{5}~;~- \dfrac{16}{5}~;0\right)$.

On calcule DI$^2 = \left(\dfrac{12}{5} - 12 \right)^2 + \left(- \dfrac{16}{5} - 0 \right)^2 + 0^2 = \left(-\dfrac{48}{5}\right)^2 + \left(- \dfrac{16}{5}\right)^2 = \dfrac{\np{2304}}{25} + \dfrac{256}{25} = \dfrac{\np{2560}}{25} = \dfrac{256 \times 10}{25}$.

Finalement DI $= \sqrt{\dfrac{256 \times 10}{25}} = \dfrac{16 \sqrt{10}}{5}$.
	\end{enumerate}
\item 





En prenant la base (ABC) rectangle en C : son aire est donc égale à :
$\dfrac{\text{AC} \times \text{BC}}{2}$.

AC$^2 = 16 + 144  + 256 = 416$ ; AC $ = \sqrt{416}$ ;

BC$^2 = 16 + 144 + 100 = 260$ ; BC $ = \sqrt{260}$.

Donc $\mathcal{A}(\text{ABC}) = \dfrac{\sqrt{416} \times \sqrt{260}}{2} = 52\sqrt{18}$

Donc avec $h = \dfrac{16 \sqrt{10}}{5}$, on obtient :

$V = \dfrac13 \times 52\sqrt{18} \times \dfrac{16 \sqrt{10}}{5} = \dfrac{\np{1664}}{3} \approx 554,6$ soit 555 à l'unité près.
\end{enumerate}

\bigskip
